% Root bilingual entry point for the Tufte-style book sample.
% Change english to chinese to switch the font mapping and localized labels.
% Compile this document with XeLaTeX.
\documentclass[nofonts,nohyper,nobib]{tufte-book}

% Allowed values: english, chinese.
\newcommand{\BookLanguage}{english}

\usepackage{Thomas-Tufte-bilingual-book}
% Font scheme:
%   tufte  = A_Tufte_Style_Book_with_VDQI font chain (Palatino + Helvetica + Bera Mono).
%   beamer = OUC-Haide-Beamer-Template font chain (RobotoSerif + Menlo + national-2),
%            loaded from ./fonts/ (a copy of OUC-Haide-Beamer-Template-main/fonts).
\SetBookFontScheme{tufte}
% \SetBookFontScheme{beamer}
\SetBookLanguage{\BookLanguage}
\OUCSetPalette{0} % 0..23, extended with the discussed thematic palettes.
% Use \textcolor{spotcolor}{...} for theme-specific point accents such as Kyoto red.

%%
% Book metadata
\title{Thomas Shang's Tufte-Style Book Template}
\date{The edition number}
\author[Thomas Shang \& Tufte-LaTeX Developers]{Thomas Shang \& Tufte-LaTeX Developers}
\publisher{Publisher of This Book}

\begin{document}

% Front matter
\frontmatter

% r.1 blank page
% \blankpage

% v.2 epigraphs
% \newpage\thispagestyle{empty}
% \openepigraph{%
% The public is more familiar with bad design than good design.
% It is, in effect, conditioned to prefer bad design, 
% because that is what it lives with. 
% The new becomes threatening, the old reassuring.
% }{Paul Rand%, {\itshape Design, Form, and Chaos}
% }
% \vfill
% \openepigraph{%
% A designer knows that he has achieved perfection 
% not when there is nothing left to add, 
% but when there is nothing left to take away.
% }{Antoine de Saint-Exup\'{e}ry}
% \vfill
% \openepigraph{%
% \ldots the designer of a new system must not only be the implementor and the first 
% large-scale user; the designer should also write the first user manual\ldots 
% If I had not participated fully in all these activities, 
% literally hundreds of improvements would never have been made, 
% because I would never have thought of them or perceived 
% why they were important.
% }{Donald E. Knuth}


% r.3 full title page
\maketitle


% v.4 copyright page
\newpage
\begin{fullwidth}
~\vfill
\thispagestyle{empty}
\setlength{\parindent}{0pt}
\setlength{\parskip}{\baselineskip}
Copyright \copyright\ \the\year\ \thanklessauthor

\par Published by \thanklesspublisher

\par tufte-latex.googlecode.com

\par Licensed under the Apache License, Version 2.0 (the ``License''); you may not
use this file except in compliance with the License. You may obtain a copy
of the License at \url{http://www.apache.org/licenses/LICENSE-2.0}. Unless
required by applicable law or agreed to in writing, software distributed
under the License is distributed on an ``AS IS'' BASIS, WITHOUT
License for the specific language governing permissions and limitations
under the License.\index{license}

\par\textit{First printing, \monthyear}

\par\medskip
This template was carefully polished and redesigned by Thomas Shang on the
basis of the upstream Tufte-style book template, refining the typography,
palette system, and bilingual workflow into the present form. Project home:
\url{https://github.com/Hatsuyuki017}.
\end{fullwidth}

% r.5 contents
\tableofcontents

% Active-palette showcase page (unnumbered, not added to ToC).
% Six 1cm colour bands of pal1..pal6 with 1.6cm vertical gaps.
% Convention (palette-independent): pal1..pal3 use black text, pal4..pal6 use white text.
% The palette's design-philosophy paragraph is shifted into the margin via \marginnote.
\cleardoublepage
\thispagestyle{empty}
\chapter*{\OUCPaletteName}
\marginnote{The factory ground tone of the template. Blue is the colour of the young scholar---from the first lecture page at dawn, through the library window at noon, to the cold lab light at midnight, blue spans the entire colour temperature of academic life. Six tones run from lightest to deepest, deliberately restrained, mirroring the arc of opening a book and closing it again. The palette does not ask to be looked at; it asks to let the words themselves be seen.}
\noindent\colorbox{oucpal1}{\parbox[c][1cm][c]{\dimexpr\linewidth-2\fboxsep\relax}{\centering\bfseries\color{black}Paper Foam\quad\textnormal{\texttt{rgb(239, 246, 255)}}}}\par
\vspace{1.6cm}
\noindent\colorbox{oucpal2}{\parbox[c][1cm][c]{\dimexpr\linewidth-2\fboxsep\relax}{\centering\bfseries\color{black}Desert Foam\quad\textnormal{\texttt{rgb(219, 234, 254)}}}}\par
\vspace{1.6cm}
\noindent\colorbox{oucpal3}{\parbox[c][1cm][c]{\dimexpr\linewidth-2\fboxsep\relax}{\centering\bfseries\color{black}Sea Smoke\quad\textnormal{\texttt{rgb(96, 165, 250)}}}}\par
\vspace{1.6cm}
\noindent\colorbox{oucpal4}{\parbox[c][1cm][c]{\dimexpr\linewidth-2\fboxsep\relax}{\centering\bfseries\color{white}Crisp Sky\quad\textnormal{\texttt{rgb(37, 99, 235)}}}}\par
\vspace{1.6cm}
\noindent\colorbox{oucpal5}{\parbox[c][1cm][c]{\dimexpr\linewidth-2\fboxsep\relax}{\centering\bfseries\color{white}Twilight Blue\quad\textnormal{\texttt{rgb(30, 58, 138)}}}}\par
\vspace{1.6cm}
\noindent\colorbox{oucpal6}{\parbox[c][1cm][c]{\dimexpr\linewidth-2\fboxsep\relax}{\centering\bfseries\color{white}Midnight Hull\quad\textnormal{\texttt{rgb(30, 41, 59)}}}}\par


\listoffigures

\listoftables

% r.7 dedication
\cleardoublepage
~\vfill
\begin{doublespace}
\noindent\fontsize{18}{22}\selectfont\itshape
\nohyphenation
Dedicated to those who appreciate \LaTeX{} 
and the work of \mbox{Edward R.~Tufte} 
and \mbox{Donald E.~Knuth}.
\end{doublespace}
\vfill
\vfill


% r.9 introduction
\cleardoublepage
\chapter*{Introduction}

This sample book discusses the design of Thomas Shang's Tufte-style book template on the basis of Edward Tufte's
books\cite{Tufte2001,Tufte1990,Tufte1997,Tufte2006}
and the use of the \doccls{tufte-book} and \doccls{tufte-handout} document classes.

\cleardoublepage

\chapter*{Palette Guide}

This sample book now exposes the same palette interface as the beamer template,
extended with the thematic palette groups discussed for this project. Select a
palette with \doccmd{OUCSetPalette}\{N\} in the preamble. The six swatches in
each entry run from \texttt{pal1} to \texttt{pal6}, ordered light to dark;
\texttt{primary} maps to \texttt{pal5}, \texttt{accent} maps to \texttt{pal6},
and the Kyoto palette additionally exposes \texttt{spotcolor} for its restrained
torii red accent.

\begin{fullwidth}
\scriptsize
\begin{multicols}{3}
\raggedcolumns
\bookpaletteguideentries
\end{multicols}
\end{fullwidth}


%%
% Start the main matter (normal chapters)

\mainmatter

\part{Template Reference}

\chapter{Showcase of Boxed Environments}
\label{ch:showcase}

This chapter mirrors the demonstration frames in the companion
\texttt{OUC-Haide-Beamer-Template} so that the book and the slide deck share
identical environment names, palette logic, and listing style.  Each section
below pairs the \emph{rendered} environment on the left with its exact
\emph{source} on the right; the active palette is \texttt{\OUCPaletteName}
and can be swapped via \doccmd{OUCSetPalette}\{N\}.

\section{Slot Semantics}\index{palette!slots}
Six palette slots drive every coloured element:
\texttt{pal1}=\hlred{paleblue} listing background,
\texttt{pal2} body fill for \docenv{Exa}/\docenv{Rmk}/\docenv{Intro},
\texttt{pal4} title fill for the same trio,
\texttt{pal5}=\hlred{primary} emphasis (theorem titles, frame titles, footer),
\texttt{pal6}=\hlred{accent} darkest swatch (definition title, chapter heads).

\section{Definition and Theorem}\index{environments!De@\texttt{De}}\index{environments!The@\texttt{The}}

\begin{fullwidth}
\begin{minipage}[t]{0.48\linewidth}
\begin{De}{Group}{de:g}
A set with an associative binary operation, an identity element, and inverses
for every element.
\end{De}
\begin{The}{Lagrange}{the:lag}
If $H\le G$ is a subgroup of a finite group, then $|H|$ divides $|G|$.
\end{The}
\end{minipage}\hfill
\begin{minipage}[t]{0.48\linewidth}
\begin{lstlisting}[language={[LaTeX]TeX}]
\begin{De}{Group}{de:g}
  A set with an associative
  operation, identity, inverses.
\end{De}
\begin{The}{Lagrange}{the:lag}
  If $H\le G$, then $|H|$
  divides $|G|$.
\end{The}
\end{lstlisting}
\end{minipage}
\end{fullwidth}

\section{Example and Remark}\index{environments!Exa@\texttt{Exa}}\index{environments!Rmk@\texttt{Rmk}}

\begin{fullwidth}
\begin{minipage}[t]{0.48\linewidth}
\begin{Exa}{}{exa:cyc}
$\mathbb{Z}/6\mathbb{Z}$ is a cyclic group of order $6$, generated by
the residue class of $1$.
\end{Exa}
\begin{Rmk}{}{rmk:tr}
The trivial group $\{e\}$ is the unique group of order $1$ up to
isomorphism.
\end{Rmk}
\end{minipage}\hfill
\begin{minipage}[t]{0.48\linewidth}
\begin{lstlisting}[language={[LaTeX]TeX}]
\begin{Exa}{}{exa:cyc}
  $\mathbb{Z}/6\mathbb{Z}$ is
  cyclic of order $6$.
\end{Exa}
\begin{Rmk}{}{rmk:tr}
  The trivial group $\{e\}$ is the
  unique group of order $1$.
\end{Rmk}
\end{lstlisting}
\end{minipage}
\end{fullwidth}

\section{Proposition, Lemma, Corollary}\index{environments!Pro@\texttt{Pro}}\index{environments!Le@\texttt{Le}}\index{environments!Co@\texttt{Co}}

\begin{fullwidth}
\begin{minipage}[t]{0.48\linewidth}
\begin{Pro}{Closure}{pro:cl}
$H\cap K$ is closed under the group operation whenever $H,K\le G$.
\end{Pro}
\begin{Le}{Cancellation}{le:can}
If $ab=ac$ in a group then $b=c$.
\end{Le}
\begin{Co}{Order divides}{co:ord}
$|\langle g\rangle|$ divides $|G|$.
\end{Co}
\end{minipage}\hfill
\begin{minipage}[t]{0.48\linewidth}
\begin{lstlisting}[language={[LaTeX]TeX}]
\begin{Pro}{Closure}{pro:cl}
  $H\cap K$ is closed.
\end{Pro}
\begin{Le}{Cancellation}{le:can}
  If $ab=ac$, then $b=c$.
\end{Le}
\begin{Co}{Order divides}{co:ord}
  $|\langle g\rangle|$
  divides $|G|$.
\end{Co}
\end{lstlisting}
\end{minipage}
\end{fullwidth}

\section{Proof and Boxed}\index{environments!Proof@\texttt{Proof}}\index{environments!Boxed@\texttt{Boxed}}

\begin{fullwidth}
\begin{minipage}[t]{0.48\linewidth}
\begin{Proof}
$x\in H\cap K \Rightarrow x\in H$ and $x\in K$, so the intersection is a
subgroup.
\end{Proof}
\begin{Boxed}
Rounded container, suited for paragraph-sized notes, derivations
or summaries that should be visually grouped together.
\end{Boxed}
\end{minipage}\hfill
\begin{minipage}[t]{0.48\linewidth}
\begin{lstlisting}[language={[LaTeX]TeX}]
\begin{Proof}
  $x\in H\cap K \Rightarrow$
  $x\in H$ and $x\in K$.
\end{Proof}
\begin{Boxed}
  Rounded container for
  paragraph-sized notes.
\end{Boxed}
\end{lstlisting}
\end{minipage}
\end{fullwidth}

\section{Intro and RightIntro}\index{environments!Intro@\texttt{Intro}}\index{environments!RightIntro@\texttt{RightIntro}}

\begin{fullwidth}
\begin{minipage}[t]{0.48\linewidth}
\begin{Intro}{Chapter Intro}{Title attached to the upper-left corner; useful as a chapter or section opener that previews the goal.}
\end{Intro}
\begin{RightIntro}{Right Intro}{Mirrored variant with the title attached to the upper-right corner.}
\end{RightIntro}
\end{minipage}\hfill
\begin{minipage}[t]{0.48\linewidth}
\begin{lstlisting}[language={[LaTeX]TeX}]
\begin{Intro}{Chapter Intro}
  {Title in the upper-left
   corner.}
\end{Intro}
\begin{RightIntro}{Right Intro}
  {Title in the upper-right
   corner.}
\end{RightIntro}
\end{lstlisting}
\end{minipage}
\end{fullwidth}

\section{UnderlineBox and Quote}\index{UnderlineBox command@\protect\hangleft{\texttt{\tuftebs}}\texttt{UnderlineBox}}\index{Quote command@\protect\hangleft{\texttt{\tuftebs}}\texttt{Quote}}

\begin{fullwidth}
\begin{minipage}[t]{0.48\linewidth}
\UnderlineBox{Underlined block bracketed by primary-coloured rules.}
\Quote{Mathematics is the queen of the sciences.}{C.\,F.~Gauss}
\end{minipage}\hfill
\begin{minipage}[t]{0.48\linewidth}
\begin{lstlisting}[language={[LaTeX]TeX}]
\UnderlineBox{Underlined block
  with two coloured rules.}
\Quote{Mathematics is the queen
  of the sciences.}{C.\,F.~Gauss}
\end{lstlisting}
\end{minipage}
\end{fullwidth}

\section{Minipage, Eq, EqL}\index{Minipage command@\protect\hangleft{\texttt{\tuftebs}}\texttt{Minipage}}\index{Eq command@\protect\hangleft{\texttt{\tuftebs}}\texttt{Eq}}\index{EqL command@\protect\hangleft{\texttt{\tuftebs}}\texttt{EqL}}

\begin{fullwidth}
\begin{minipage}[t]{0.48\linewidth}
\Minipage{0.46}{\textbf{Left.}\\Side-by-side layout helper.}{0.46}{\textbf{Right.}\\Widths must sum below $1$.}
\Eq{a^2 + b^2 = c^2}
\EqL{e^{i\pi} + 1 = 0}{eq:euler}
\end{minipage}\hfill
\begin{minipage}[t]{0.48\linewidth}
\begin{lstlisting}[language={[LaTeX]TeX}]
\Minipage{0.46}{Left content}
         {0.46}{Right content}
\Eq{a^2 + b^2 = c^2}
\EqL{e^{i\pi} + 1 = 0}{eq:euler}
\end{lstlisting}
\end{minipage}
\end{fullwidth}

\section{Info, Warn, Result, Goal Boxes}\index{environments!infobox@\texttt{infobox}}\index{environments!warnbox@\texttt{warnbox}}\index{environments!resultbox@\texttt{resultbox}}\index{environments!goalbox@\texttt{goalbox}}

\begin{fullwidth}
\begin{minipage}[t]{0.48\linewidth}
\begin{infobox}
Background, definitions, and context.
\end{infobox}
\begin{warnbox}
Caveats, assumptions, and risks.
\end{warnbox}
\begin{resultbox}
Findings, conclusions, and final outputs.
\end{resultbox}
\begin{goalbox}
\centering A short highlighted message.
\end{goalbox}
\end{minipage}\hfill
\begin{minipage}[t]{0.48\linewidth}
\begin{lstlisting}[language={[LaTeX]TeX}]
\begin{infobox}...\end{infobox}
\begin{warnbox}...\end{warnbox}
\begin{resultbox}...\end{resultbox}
\begin{goalbox}...\end{goalbox}
\end{lstlisting}
\end{minipage}
\end{fullwidth}

\section{Listing Style for Code}\index{lstlisting style@\texttt{oucbooklisting} style}

The \texttt{oucbooklisting} style is registered as the default for the
\docpkg{listings} package and follows the palette: paleblue background,
primary keywords, accent strings, grey frame.  The same style is shared with
the beamer template.

\subsection{Matlab}
\begin{lstlisting}[language=Matlab]
function y = laplacian(u, dx)
% Five-point discrete Laplacian on a uniform grid.
y = (circshift(u,[0 -1]) + circshift(u,[0 1]) ...
   + circshift(u,[-1 0]) + circshift(u,[ 1 0]) - 4*u) / dx^2;
end
\end{lstlisting}

\subsection{Python}
\begin{lstlisting}[language=Python]
def softmax(x):
    """Numerically stable softmax."""
    z = x - x.max(axis=-1, keepdims=True)
    e = np.exp(z)
    return e / e.sum(axis=-1, keepdims=True)
\end{lstlisting}

\subsection{Shell}
\begin{lstlisting}[language=bash]
latexmk -xelatex -interaction=nonstopmode Thomas-Tufte-book.tex
makeindex Thomas-Tufte-book.idx
latexmk -xelatex -interaction=nonstopmode Thomas-Tufte-book.tex
\end{lstlisting}

\subsection{LaTeX}
\begin{lstlisting}[language={[LaTeX]TeX}]
\usepackage{Thomas-Tufte-bilingual-book}
\SetBookLanguage{english}
\OUCSetPalette{0}
\begin{document}
  \maketitle
  \tableofcontents
  \mainmatter
\end{document}
\end{lstlisting}

\cleardoublepage


\part{The First Part}
\input{chapter-en/05-the-design-of-tufte-s-books}
\input{chapter-en/06-on-the-use-of-the-tufte-book}
\input{chapter-en/07-customizing}
\input{chapter-en/08-compatibility-issues}
\input{chapter-en/09-troubleshooting-and-support}







\backmatter

% \bibliography{sample-handout}
% \bibliographystyle{plainnat}


\chapter*{Acknowledgements}
\addcontentsline{toc}{chapter}{Acknowledgements}

This template stands on the shoulders of several open works. The author thanks:

\begin{itemize}
  \item The \textbf{Tufte-LaTeX} project, whose \texttt{tufte-book} class is the typographic
    foundation of this template.\\
    \url{https://github.com/Tufte-LaTeX/tufte-latex}
  \item The \textbf{OUC Haide Beamer Template}, whose palette mechanism and visual
    language this book template inherits and extends.\\
    \url{https://github.com/Hatsuyuki017/OUC-Haide-Beamer-Template}
  \item \textbf{Mattia Puddu}, whose Tufte-style book layout with VDQI-inspired title and
    contents pages directly inspired the front-matter design of this template.\\
    \url{https://mattiapuddu25.github.io/index.html}
\end{itemize}



\printindex

% Palette Introduction --- English.
% Auto-generated by gen_palette_intro.py.  Do not edit by hand.
\chapter{Palette Introduction}\label{ch:palette-intro}

This chapter walks through all twenty-four themed palettes in order. Every entry keeps a fixed three-part structure---a \emph{Colour Philosophy} stating the worldview behind the palette, a \emph{six-colour table} listing the values and gloss of \texttt{pal1} through \texttt{pal6}, and a closing subsection (\emph{Design Notes}, \emph{Colour Structure}, or \emph{Aesthetic Rules} as appropriate) unfolding either per-colour annotations or structural principles. Palettes are laid out sequentially with natural page breaks; the only constraint enforced is that each six-colour table is kept intact on a single page.

% ---------------------------------------------------------------------------
\bookpalettecardhead{0\enspace OUC Default}
\palettesubsec{Colour Philosophy}
\noindent This is the factory ground tone of the template. Blue is the colour of the young scholar---from the first lecture page at dawn, through the library window at noon, to the cold lab light at midnight, blue spans the entire colour temperature of academic life. Six tones run from lightest to deepest, deliberately restrained, mirroring the arc of opening a book and closing it again. The palette does not ask to be looked at; it asks to let the words themselves be seen. This is the quietest possible contract between writer and reader, and the way the OUC Haide LaTeX template hopes to be used for the long haul.

\begin{fullwidth}\small\noindent
\begin{tabularx}{\linewidth}{@{}clllX@{}}
\toprule\rowcolor{booktableheadcolor}
\booktableheadcell{} & \booktableheadcell{Name} & \booktableheadcell{Original} & \booktableheadcell{RGB} & \booktableheadcell{Gloss} \\
\midrule
\bookcardswatch{239,246,255} & Paper Foam & Paper Foam & 239, 246, 255 & The faintest blue a page reflects at the moment of opening, almost white \\
\bookcardswatch{219,234,254} & Desert Foam & Desert Foam & 219, 234, 254 & A page in shadow under cold lab light, clean and weightless \\
\bookcardswatch{96,165,250} & Sea Smoke & Sea Smoke & 96, 165, 250 & The window sky, the breath a scholar takes between paragraphs \\
\bookcardswatch{37,99,235} & Crisp Sky & Crisp Sky & 37, 99, 235 & A chalk stroke on the board, the colour of argument \\
\bookcardswatch{30,58,138} & Twilight Blue & Twilight Blue & 30, 58, 138 & The deepest blue of contemplation, near the edge of dream \\
\bookcardswatch{30,41,59} & Midnight Hull & Midnight Hull & 30, 41, 59 & The final line a pen leaves in the notebook \\
\bottomrule
\end{tabularx}\end{fullwidth}

\palettesubsec{Colour Mapping of the Reading Arc}
\begin{palettebullets}
\item \textbf{Paper Foam} --- Paper Foam is the breath of "opening the book"---near-white, allowing any text to enter the reader's eye with the smallest visual cost
\item \textbf{Frost Crystal} --- Frost Crystal carries the second-lightest weight, replacing pure white over long reads to lower fatigue without drawing attention
\item \textbf{Sky Glass} --- Sky Glass is the only mid-tone with a hint of nature, pulling the page out of pure artificial colour and back to a glance out the window
\item \textbf{Lecture Sapphire} --- Lecture Sapphire bears the duty of emphasis, the only tone permitted to be "loud"---yet loud with discipline, like chalk landing a stress on the board
\item \textbf{Reverie Navy} --- Reverie Navy presses thought to its deepest layer, suiting long quotations, pivotal theorems, and chapter openings
\item \textbf{Slate Ink} --- Slate Ink is almost black yet keeps a blue heart---the colour every line of body text and every structural rule finally lands in; the book closes inside it
\end{palettebullets}
\par\bigskip
% ---------------------------------------------------------------------------
\bookpalettecardhead{1\enspace Brunneophobia}
\palettesubsec{Colour Philosophy}
\noindent The name literally reads "fear of brown", yet the palette faces that prejudice with nothing but brown. Brown is neither cheap nor dated; it is the warmth that modern taste has deliberately pushed aside. Six tones descend from desert sand and smoky tangerine through rust to the deepest scorched earth, like walking into a wood-panelled studio of old books, the air thick with leather, tobacco, and damp timber. Courage in colour is not brilliance---it is daring to be the shade everyone avoids, and brown, returned to itself, becomes the very definition of grown-up dignity.

\begin{fullwidth}\small\noindent
\begin{tabularx}{\linewidth}{@{}clllX@{}}
\toprule\rowcolor{booktableheadcolor}
\booktableheadcell{} & \booktableheadcell{Name} & \booktableheadcell{Original} & \booktableheadcell{RGB} & \booktableheadcell{Gloss} \\
\midrule
\bookcardswatch{238,211,180} & Lighthouse & Lighthouse & 238, 211, 180 & The faintest warm white of desert dawn, the palette\'s breath \\
\bookcardswatch{213,148,79} & Sand Lily & Sand Lily & 213, 148, 79 & Tangerine seen through smoke, warm but never loud \\
\bookcardswatch{213,148,79} & Reef Coral & Reef Coral & 213, 148, 79 & The same hue repeated---a deliberate narrowing of chroma \\
\bookcardswatch{180,69,15} & Tidal Mauve & Tidal Mauve & 180, 69, 15 & Iron left in damp air for decades, the colour of blood and soil \\
\bookcardswatch{86,67,53} & Inkwell & Inkwell & 86, 67, 53 & The shade of an old sofa and a leather-bound spine, dense with years \\
\bookcardswatch{42,23,14} & Deep Tide & Deep Tide & 42, 23, 14 & The deepest brown left in the ash after the campfire dies \\
\bottomrule
\end{tabularx}\end{fullwidth}

\palettesubsec{The Internal Order of the Brown Family}
\begin{palettebullets}
\item \textbf{Desert Foam} --- Desert Foam, as the lightest tone, carries every duty of breathing and negative space---the only brown in the set that does not speak
\item \textbf{Smoke Tangerine} --- Smoke Tangerine governs the warmth of the middle band, returning the palette from gravity to daily life---the friendliest interface with the reader
\item \textbf{Beeswax Honey} --- Beeswax Honey shares a hue with Smoke Tangerine---a deliberate clenching, an argument that restraint itself is force
\item \textbf{Rust Crimson} --- Rust Crimson is the edge of the palette---it pushes gentle browns into an aggressive red, the only tone permitted to wound
\item \textbf{Aged Leather} --- Aged Leather supplies the texture of archives---the deep tone shared by libraries, dossiers, and leather bindings, locking the physicality of reading into colour
\item \textbf{Scorched Earth} --- Scorched Earth is the root of structure; every warm tone above it stands because it does---the final step in overcoming the fear of brown is to admit that black, too, is brown
\end{palettebullets}
\par\bigskip
% ---------------------------------------------------------------------------
\bookpalettecardhead{2\enspace Van Dyke}
\palettesubsec{Colour Philosophy}
\noindent In tribute to Anthony van Dyck, the seventeenth-century Flemish portraitist whose ground colour, Van Dyke Brown, layered glaze upon glaze to lend aristocratic skin its translucent glow. The palette grows out of his canvas itself: powdered flesh tones, lilac silk, crimson velvet, deep brown ground---all standard materials of Baroque portraiture. There is no drama between the six tones, only the composure left after fading; together they do not form a painting so much as the atmosphere a painting leaves on a wall once time has washed across it.

\begin{fullwidth}\small\noindent
\begin{tabularx}{\linewidth}{@{}clllX@{}}
\toprule\rowcolor{booktableheadcolor}
\booktableheadcell{} & \booktableheadcell{Name} & \booktableheadcell{Original} & \booktableheadcell{RGB} & \booktableheadcell{Gloss} \\
\midrule
\bookcardswatch{236,194,188} & Carnation Skin & Vleeskleur & 236, 194, 188 & The thin-glazed pink of an aristocrat\'s cheek in Baroque portraiture \\
\bookcardswatch{169,159,191} & Silken Wisteria & Zijden Sering & 169, 159, 191 & The lilac silk catches in shadow, elegant and cool \\
\bookcardswatch{169,159,191} & Smoke Lilac & Rooklila & 169, 159, 191 & As above---both sides of the silk are barely distinct on canvas \\
\bookcardswatch{191,113,133} & Rose Velvet & Rozenfluweel & 191, 113, 133 & The rose that bleeds through velvet, deep without being loud \\
\bookcardswatch{68,60,94} & Plum Plush & Pruimenpurper & 68, 60, 94 & The deep plum of the backdrop drapery, throwing the highlights brighter \\
\bookcardswatch{61,43,39} & Van Dyke Brown & Van Dijck-bruin & 61, 43, 39 & Van Dyck\'s signature ground, mixed from bitumen pigment \\
\bottomrule
\end{tabularx}\end{fullwidth}

\palettesubsec{Colour Mapping of Portrait Depth}
\begin{palettebullets}
\item \textbf{Carnation Skin} --- Carnation Skin sits at the front---the cheek, the hand, the neck of the figure, the first "person" the viewer meets
\item \textbf{Silken Wisteria} --- Silken Wisteria is the silk's highlight, holding softness of refraction in the hue---elegant with distance
\item \textbf{Smoke Lilac} --- Smoke Lilac shares its hue with Silken Wisteria, a deliberate repetition that echoes Baroque painting's love of monochromatic layering
\item \textbf{Rose Velvet} --- Rose Velvet is the mid-ground gown, lifting and pushing the figure forward---the colour-temperature pivot of the set
\item \textbf{Plum Plush} --- Plum Plush is the curtain behind, lifting the figure out of darkness while almost escaping notice itself
\item \textbf{Van Dyke Brown} --- Van Dyke Brown is the "ground" of the whole work---the colour deposited by pigment, canvas, frame, and wall together; every highlight stands only because it does
\end{palettebullets}
\par\bigskip
% ---------------------------------------------------------------------------
\bookpalettecardhead{3\enspace Back in Black}
\palettesubsec{Colour Philosophy}
\noindent In tribute to AC/DC's 1980 album, though the palette itself is more restrained than rock. Pink and grey soften the aggression of black, pushing the set towards a post-punk aesthetic: leather jacket, smoke, the residual warmth right after the stage lights die. Pink is the scar, black is the night---every grey between them is the echo youth leaves inside the body. This is a grown-up version of rebellion, with its edge folded inside a composed surface.

\begin{fullwidth}\small\noindent
\begin{tabularx}{\linewidth}{@{}clllX@{}}
\toprule\rowcolor{booktableheadcolor}
\booktableheadcell{} & \booktableheadcell{Name} & \booktableheadcell{Original} & \booktableheadcell{RGB} & \booktableheadcell{Gloss} \\
\midrule
\bookcardswatch{240,217,228} & Powder Pink & Powder Pink & 240, 217, 228 & The only undefended colour in the dark, the opposite of rebellion \\
\bookcardswatch{193,160,172} & Dust Rose & Dust Rose & 193, 160, 172 & Pink with a film of smoke, dignity after the brightness has gone \\
\bookcardswatch{193,160,172} & Ember Pink & Ember Pink & 193, 160, 172 & Repeated to strengthen the mid-band; two semantics of pink, side by side \\
\bookcardswatch{128,108,121} & Smoky Mauve & Smoky Mauve & 128, 108, 121 & The smoke-grey purple under stage light---neither bright nor dim \\
\bookcardswatch{74,63,75} & Black Velvet & Black Velvet & 74, 63, 75 & The shadow of a leather jacket, between black and purple \\
\bookcardswatch{22,19,21} & Inkwell Black & Inkwell Black & 22, 19, 21 & The deepest anchor of the set, all but pure black \\
\bottomrule
\end{tabularx}\end{fullwidth}

\palettesubsec{Colour Mapping of a Post-Punk Aesthetic}
\begin{palettebullets}
\item \textbf{Powder Pink} --- Powder Pink is the unguarded entry---the only soft door above the black ground; rebellion's tender opposite is rebellion's core
\item \textbf{Dust Rose} --- Dust Rose presses the brightness out of pink---the after-tone of youth, elegantly admitting it is no longer young
\item \textbf{Ember Pink} --- Ember Pink shares its hue with Dust Rose---a repeat-as-emphasis figure that promotes the mid-band into the true centre of the palette
\item \textbf{Smoky Mauve} --- Smoky Mauve is the smoke under the lights, bridging pink and black so the two extremes look like two faces of one mood
\item \textbf{Black Velvet} --- Black Velvet is the shadow of a leather jacket, breathing with purple---black does not have to be dead; it can be a substance worn on the body
\item \textbf{Inkwell Black} --- Inkwell Black is the deepest anchor, near pure black yet keeping a violet undertone---the period of the palette, though even after a period the next song still plays
\end{palettebullets}
\par\bigskip
% ---------------------------------------------------------------------------
\bookpalettecardhead{4\enspace Belle of the Ball}
\palettesubsec{Colour Philosophy}
\noindent The most-watched woman at the ball never wears primaries---she wears the old tones that countless candle flames have warmed: terracotta, verdigris, olive, deep crimson. The palette refuses twentieth-century synthetic dyes and is drawn instead from the mixed aesthetic of the nineteenth-century European salon and its colonial tropics---fireplace, mahogany, evening curtain, wine-stained linen. Being looked at depends not on brilliance but on that "lived-through" elegance every tone wears once time has softened it.

\begin{fullwidth}\small\noindent
\begin{tabularx}{\linewidth}{@{}clllX@{}}
\toprule\rowcolor{booktableheadcolor}
\booktableheadcell{} & \booktableheadcell{Name} & \booktableheadcell{Original} & \booktableheadcell{RGB} & \booktableheadcell{Gloss} \\
\midrule
\bookcardswatch{226,203,192} & Champagne Pearl & Champagne Pearl & 226, 203, 192 & The pearled skin tone of a woman under candlelight in the ballroom \\
\bookcardswatch{206,171,150} & Rose Terracotta & Rose Terracotta & 206, 171, 150 & A mid-tone between vintage lipstick and terracotta---the warmth of being watched \\
\bookcardswatch{210,135,106} & Glazed Clay & Glazed Clay & 210, 135, 106 & The orange-red of glazed pottery---warm, mature, hand-thrown \\
\bookcardswatch{229,74,57} & Chilli Scarlet & Chilli Scarlet & 229, 74, 57 & The scarlet of chilli peppers drying in tropical sun---the bright accent of the set \\
\bookcardswatch{118,118,44} & Olive Verdigris & Olive Verdigris & 118, 118, 44 & Olive leaves oxidised to a neutral green with a bronze edge \\
\bookcardswatch{53,77,4} & Moss Deepwood & Moss Deepwood & 53, 77, 4 & The forest-dark tone of mahogany furniture and deep curtains \\
\bottomrule
\end{tabularx}\end{fullwidth}

\palettesubsec{Colour Strategy of the One Being Watched}
\begin{palettebullets}
\item \textbf{Champagne Pearl} --- Champagne Pearl is skin warmed by candlelight---the breath of the palette and the most undefended layer of the watched figure
\item \textbf{Rose Terracotta} --- Rose Terracotta carries vintage allure---a lipstick filtered through time, both ornament and proof of standing
\item \textbf{Glazed Clay} --- Glazed Clay is the colour of handcraft, banishing industry from the palette and keeping in its place the warmth of a palm and the substance of earth
\item \textbf{Chilli Scarlet} --- Chilli Scarlet is the only tone permitted to be bright---the red of a hem, a ring, a final lipstick touch---and only ever in small area
\item \textbf{Olive Verdigris} --- Olive Verdigris is the cool pivot, pulling every tropical warm back into the restraint of a European salon---the palette's thermostat
\item \textbf{Moss Deepwood} --- Moss Deepwood is the background and the furniture, the gravity of the set---without it, every bright tone would lose its pivot
\end{palettebullets}
\par\bigskip
% ---------------------------------------------------------------------------
\bookpalettecardhead{5\enspace Pine Tree}
\palettesubsec{Colour Philosophy}
\noindent Anchored on a single northern pine, the palette descends from canopy to root: warm gold where sun strikes the new needles, the loosened old needles tossed by wind, rose-clay shadow under the afternoon canopy, ochre seeping from cracks in the trunk, and the dark forest ink buried in the humus below. The set deliberately leaves "green" behind---a real pine, seen carefully, is far more complex than green; its colour is a chronology jointly written by sunlight, weather, soil, wounds, and shadow.

\begin{fullwidth}\small\noindent
\begin{tabularx}{\linewidth}{@{}clllX@{}}
\toprule\rowcolor{booktableheadcolor}
\booktableheadcell{} & \booktableheadcell{Name} & \booktableheadcell{Original} & \booktableheadcell{RGB} & \booktableheadcell{Gloss} \\
\midrule
\bookcardswatch{238,200,111} & Sunlit Needle & Sunlit Needle & 238, 200, 111 & Warm gold where sunlight strikes the new needles \\
\bookcardswatch{222,166,32} & Honey Resin & Honey Resin & 222, 166, 32 & Pine resin oozing from a crack in the bark, the tree\'s self-healing gold \\
\bookcardswatch{222,166,32} & Aged Gold & Aged Gold & 222, 166, 32 & The same gold returned, steadying the mid-band \\
\bookcardswatch{177,120,133} & Shadow Rose & Shadow Rose & 177, 120, 133 & Rose-clay in the afternoon shade, with a violet edge \\
\bookcardswatch{167,88,26} & Bark Ochre & Bark Ochre & 167, 88, 26 & The ochre brown of the trunk itself---the tree\'s most honest colour \\
\bookcardswatch{43,47,34} & Forest Ink & Forest Ink & 43, 47, 34 & The dark forest ink of humus and root \\
\bottomrule
\end{tabularx}\end{fullwidth}

\palettesubsec{Colour Mapping of a Pine Tree's Chronology}
\begin{palettebullets}
\item \textbf{Sunlit Needle} --- Sunlit Needle is the pine tree's "today"---the newly-grown, sunlight-chosen tip that has no story yet
\item \textbf{Honey Resin} --- Honey Resin is the wound---resin sealing a crack into amber, the tone most charged with time in the entire palette
\item \textbf{Aged Gold} --- Aged Gold shares its hue with Honey Resin---a "say it again" of the mid-band gold, time recited once more in colour
\item \textbf{Shadow Rose} --- Shadow Rose is the violet shadow the pine casts on the ground---the breathing interface between tree and earth
\item \textbf{Bark Ochre} --- Bark Ochre is the trunk itself, the colour needing no explanation---a pine's very substance, named
\item \textbf{Forest Ink} --- Forest Ink is the root colour, sunk in humus, near black yet still green---reminding us that a living root is never quite the same as black
\end{palettebullets}
\par\bigskip
% ---------------------------------------------------------------------------
\bookpalettecardhead{6\enspace Provence Blue}
\palettesubsec{Colour Philosophy}
\noindent Provençal blue does not belong to the sea---it belongs to the blue-grey beyond a lavender field that the setting sun has desaturated, to the peacock paint flaking from a limestone window-frame after years, to the warm air shared by old linen sheets drying in a courtyard and a rusted iron handle. The palette refuses the postcard-Mediterranean romance and chooses an inland, rural, time-softened southern-French blue---its romance lives not in the brightness of the colour but in the trace the colour leaves once it has been used.

\begin{fullwidth}\small\noindent
\begin{tabularx}{\linewidth}{@{}clllX@{}}
\toprule\rowcolor{booktableheadcolor}
\booktableheadcell{} & \booktableheadcell{Name} & \booktableheadcell{Original} & \booktableheadcell{RGB} & \booktableheadcell{Gloss} \\
\midrule
\bookcardswatch{170,188,175} & Olive Silver & Argent d'Olivier & 170, 188, 175 & The silver-grey underside of an olive leaf, with a green note \\
\bookcardswatch{137,156,154} & Mist Mirror & Miroir de Brume & 137, 156, 154 & Morning mist pressed onto water, a grey-blue \\
\bookcardswatch{137,156,154} & Spent Rain & Pluie Pass\'ee & 137, 156, 154 & Repeated, steadying the middle band \\
\bookcardswatch{110,124,139} & Window Stone Blue & Bleu de Volet & 110, 124, 139 & The peacock blue of a limestone window-frame, faded \\
\bookcardswatch{82,92,121} & Lavender Dusk & Cr\'epuscule Lavande & 82, 92, 121 & The deep violet-blue at the far edge of a lavender field under sunset \\
\bookcardswatch{53,66,94} & Olive Night & Nuit d'Olivier & 53, 66, 94 & The deep blue night sky above the olive grove, with the smell of earth \\
\bottomrule
\end{tabularx}\end{fullwidth}

\palettesubsec{Colour Structure of Rural Southern-French Blue}
\begin{palettebullets}
\item \textbf{Olive Silver} --- Olive Silver is the only green in the palette---the silver flash of an olive leaf turning in the wind, Provençal air made matter
\item \textbf{Mist Mirror} --- Mist Mirror is the dawn reflection on water, lending blue from the sky to the ground---the most breathable mid-tone
\item \textbf{Spent Rain} --- Spent Rain shares its hue with Mist Mirror, the colour temperature of the second day after a rain---repeated to deepen the humidity of the palette
\item \textbf{Window Stone Blue} --- Window Stone Blue pulls the palette out of nature back into the made---a hand-trace painted many times and faded many times more
\item \textbf{Lavender Dusk} --- Lavender Dusk is the soul of the set---it pushes lavender from "violet" into "deep violet-blue", and that single step is exactly what makes Provence Provence
\item \textbf{Olive Night} --- Olive Night is the root, gathering every daytime blue back into a single sky, near black yet still blue
\end{palettebullets}
\par\bigskip
% ---------------------------------------------------------------------------
\bookpalettecardhead{7\enspace Fresco Blue}
\palettesubsec{Colour Philosophy}
\noindent From the church-fresco skies of the Italian Renaissance---the azure that fresco technique drew into the lime plaster and that five centuries of candle smoke have since oxidised. It is not the real colour of sky but the "sacred sky" mixed by craftsmen out of lapis lazuli, malachite, and lime: a blue between earth and heaven. The palette steps from the palest sky to the deepest night, mimicking the way fresco blue deepens as the wall falls from dome to floor---a sky housed inside architecture.

\begin{fullwidth}\small\noindent
\begin{tabularx}{\linewidth}{@{}clllX@{}}
\toprule\rowcolor{booktableheadcolor}
\booktableheadcell{} & \booktableheadcell{Name} & \booktableheadcell{Original} & \booktableheadcell{RGB} & \booktableheadcell{Gloss} \\
\midrule
\bookcardswatch{166,224,244} & Vault Sky & Cielo di Volta & 166, 224, 244 & The palest blue at the highest point of a vaulted fresco \\
\bookcardswatch{71,169,207} & Cobalt Glaze & Vetrina Cobalto & 71, 169, 207 & The mid-blue mixed from lapis lazuli and lime wash \\
\bookcardswatch{71,169,207} & Lapis Renewed & Lapislazzuli & 71, 169, 207 & Repeated, steadying the mid-band of the fresco wall \\
\bookcardswatch{9,121,158} & Apse Blue & Azzurro d'Abside & 9, 121, 158 & The deep blue still unfaded after five centuries in the apse \\
\bookcardswatch{4,75,102} & Vigil Blue & Azzurro della Veglia & 4, 75, 102 & The dark blue cast onto the lower fresco during night vigil \\
\bookcardswatch{2,31,46} & Crypt Blue & Azzurro della Cripta & 2, 31, 46 & The deepest blue, almost vanished, in crypt and shadowed gallery \\
\bottomrule
\end{tabularx}\end{fullwidth}

\palettesubsec{The Vertical Spectrum of a Fresco Sky}
\begin{palettebullets}
\item \textbf{Vault Sky} --- Vault Sky is the highest point of the fresco, the pale blue light bounces off the plaster---the "heaven" of the palette
\item \textbf{Cobalt Glaze} --- Cobalt Glaze is the dominant mid-tone, the full blue a worshipper sees when they look up inside the church
\item \textbf{Lapis Renewed} --- Lapis Renewed shares its hue with Cobalt Glaze, the "reaffirmation"---fresco blue never relies on one stroke; it relies on layering upon layering
\item \textbf{Apse Blue} --- Apse Blue lives in the deepest part of the church where light barely reaches, the fresco there congealed into an unfading depth
\item \textbf{Vigil Blue} --- Vigil Blue is the night colour-temperature---the same wall, under oil lamp, reveals a dark blue utterly unlike its daytime self
\item \textbf{Crypt Blue} --- Crypt Blue is where the fresco is nearly swallowed by darkness, near black yet still blue---the most silent act of "believing" in the palette
\end{palettebullets}
\par\bigskip
% ---------------------------------------------------------------------------
\bookpalettecardhead{8\enspace Monet}
\palettesubsec{Colour Philosophy}
\noindent The central proposition of Impressionism is that light itself has colour. Monet painted the same haystack, the same cathedral, the same lily pond again and again across a lifetime, all to catch the fact that "the same object owns different colours at different hours". The palette reproduces no single Monet, but extracts from the body of his work six hues most diagnostic of Impressionist identity: creamy canvas, rose-touched hay, green meadow, summer pond, morning-mist blue, deep reed shadow---their relation is not composition but the refraction of one beam of light across different substances.

\begin{fullwidth}\small\noindent
\begin{tabularx}{\linewidth}{@{}clllX@{}}
\toprule\rowcolor{booktableheadcolor}
\booktableheadcell{} & \booktableheadcell{Name} & \booktableheadcell{Original} & \booktableheadcell{RGB} & \booktableheadcell{Gloss} \\
\midrule
\bookcardswatch{247,244,213} & Canvas Cream & Toile Cr\`eme & 247, 244, 213 & The creamy ground of an untouched Monet canvas \\
\bookcardswatch{211,150,140} & Hay Rose & Foin Rose & 211, 150, 140 & The rose tipping the haystack in Monet\'s evening series \\
\bookcardswatch{211,150,140} & Evening Glow & Lueur du Soir & 211, 150, 140 & Repeated, prolonging the colour-temperature of dusk \\
\bookcardswatch{131,153,88} & Meadow Green & Pr\'e Vert & 131, 153, 88 & The mid-green of the Giverny meadow in the afternoon \\
\bookcardswatch{16,86,102} & Pond Aqua & \'Etang Aigue & 16, 86, 102 & The deep blue mirrored beneath the water in the lily series \\
\bookcardswatch{10,51,35} & Reed Shadow & Ombre des Roseaux & 10, 51, 35 & The deep green at the base of the pondside reeds, near ink \\
\bottomrule
\end{tabularx}\end{fullwidth}

\palettesubsec{Colour Structure of Impressionist Light}
\begin{palettebullets}
\item \textbf{Canvas Cream} --- Canvas Cream is the palette's "unfinished"---the breath of the canvas itself before any of Monet's brush strokes
\item \textbf{Hay Rose} --- Hay Rose is the most diagnostic Impressionist stroke---it pushes "yellow" into "rose" precisely because evening light has rewritten every local colour
\item \textbf{Evening Glow} --- Evening Glow shares the hue of Hay Rose, the "prolonged dusk"---Impressionism returned to the same subject again and again precisely to catch this line of colour
\item \textbf{Meadow Green} --- Meadow Green is the mid-tone of the palette, the true green of sunlight on new grass, neither cool nor warm
\item \textbf{Pond Aqua} --- Pond Aqua is the most diagnostic hue of late Monet's Water Lilies, lending blue from sky to water---the emotional pivot of the palette
\item \textbf{Reed Shadow} --- Reed Shadow is the root, nearly ink yet still green---Impressionism's shadows never truly accept black
\end{palettebullets}
\par\bigskip
% ---------------------------------------------------------------------------
\bookpalettecardhead{9\enspace Narcissus}
\palettesubsec{Colour Philosophy}
\noindent Not the Narcissus of Greek myth, but the botanical narcissus---a flower drawn from a bulb that opens at the close of winter. Its colour line runs from the pallid scale of the bulb to the gold of the corona and on to the deep brown of root in soil: a material narrative from underground to surface. The palette deliberately removes "green leaf", the most obvious element, so that the eye can focus on the flower's full arc---gestation, opening, return to earth---a botanical spectrum structured along the vertical axis of time.

\begin{fullwidth}\small\noindent
\begin{tabularx}{\linewidth}{@{}clllX@{}}
\toprule\rowcolor{booktableheadcolor}
\booktableheadcell{} & \booktableheadcell{Name} & \booktableheadcell{Original} & \booktableheadcell{RGB} & \booktableheadcell{Gloss} \\
\midrule
\bookcardswatch{221,213,200} & Bulb Milk & Bulbus Lacteus & 221, 213, 200 & The milk-white of a freshly cut narcissus bulb in cross-section \\
\bookcardswatch{185,149,144} & Petal Brown-Rose & Petalum Roseobrunneum & 185, 149, 144 & A fallen petal pressed into soil, dust rose \\
\bookcardswatch{185,149,144} & Faded Rose & Rosa Defloruit & 185, 149, 144 & The same colour returned---the "use it again" gesture \\
\bookcardswatch{199,149,72} & Narcissus Gold & Aurum Narcissi & 199, 149, 72 & The warm gold at the heart of the corona, the palette\'s signal colour \\
\bookcardswatch{190,108,26} & Honey Bronze & Aes Mellitum & 190, 108, 26 & The orange-brown of the bulb\'s husk, sitting between soil and flower \\
\bookcardswatch{110,60,31} & Soil Deep & Humus Profundus & 110, 60, 31 & The humus dark in which the narcissus roots, the colour of root \\
\bottomrule
\end{tabularx}\end{fullwidth}

\palettesubsec{Colour Slices of a Narcissus Life Cycle}
\begin{palettebullets}
\item \textbf{Bulb Milk} --- Bulb Milk is the narcissus before it opens---all the latent colour a plant already owns before it becomes flower
\item \textbf{Petal Brown-Rose} --- Petal Brown-Rose is the petal after the fall, the colour rewritten by soil and rain once it leaves the stem
\item \textbf{Faded Rose} --- Faded Rose shares its hue with Petal Brown-Rose---"the same fallen petal, remembered"---one falling, returning a second time in memory
\item \textbf{Narcissus Gold} --- Narcissus Gold is the core of the palette, the warm gold at the heart of the corona that appears only in early spring---the colour of "opening" itself
\item \textbf{Honey Bronze} --- Honey Bronze pushes the flower's gold towards the earth's ochre---the transition from open flower to fallen petal, from air back to ground
\item \textbf{Soil Deep} --- Soil Deep is where the narcissus comes from and returns to, near black yet warm---the palette's closing argument that every blooming arrives from darkness and returns to it
\end{palettebullets}
\par\bigskip
% ---------------------------------------------------------------------------
\bookpalettecardhead{10\enspace Roman Empire}
\palettesubsec{Colour Philosophy}
\noindent A palette indexed by the material legacy of the Roman Empire: marble, legionary cloak, senatorial toga, laurel crown, triumphal gold, murex purple. The six tones are not designed but grown from the city of Rome itself---the stone of temples, the blood of battlefields, the authority of the forum, the mark of emperors---together forming a political archaeology of colour. Every shade has a traceable material origin.

\begin{fullwidth}\small\noindent
\begin{tabularx}{\linewidth}{@{}clllX@{}}
\toprule\rowcolor{booktableheadcolor}
\booktableheadcell{} & \booktableheadcell{Name} & \booktableheadcell{Original} & \booktableheadcell{RGB} & \booktableheadcell{Gloss} \\
\midrule
\bookcardswatch{236,232,225} & Carrara Marble & Marmor Carrariense & 236, 232, 225 & The cool, faintly warm white of Carrara marble---the matter of temples and statues \\
\bookcardswatch{212,175,55} & Gloria Aurum & Aurum Gloriae & 212, 175, 55 & A neutral gold of triumphal coinage, halfway between mint and idol \\
\bookcardswatch{74,110,65} & Laurel Viridis & Laurus Viridis & 74, 110, 65 & The matte deep green of the laurel crown, taken from dried leaves \\
\bookcardswatch{180,30,30} & Legion Crimson & Sagum Legionarium & 180, 30, 30 & The saturated red of the legionary sagum---iron, and blood \\
\bookcardswatch{120,20,40} & Senate Bordeaux & Clavus Senatorius & 120, 20, 40 & The deeper wine of the senatorial toga\'s clavus, restraint over force \\
\bookcardswatch{88,28,90} & Tyrian Purple & Purpura Tyria & 88, 28, 90 & Murex dye, worth its weight in gold, reserved for emperors \\
\bottomrule
\end{tabularx}\end{fullwidth}

\palettesubsec{Design Notes}
\begin{palettebullets}
\item \textbf{Carrara Marble} --- Drawn from Carrara marble itself, a cool white with a faint warm grey, evoking the substance of temple and statue
\item \textbf{Legion Crimson} --- The vivid battle-red of the legionary Sagum, saturated and forceful---a symbol of iron and blood
\item \textbf{Senate Bordeaux} --- The deep wine red of the Clavus stripe on the senatorial toga, more restrained, more aristocratic than the legion's red
\item \textbf{Laurel Viridis} --- The matte deep green of laurel leaves---not emerald, but the steady tone of the leaf once dried
\item \textbf{Gloria Aurum} --- A neutral gold of triumphal procession, halfway between minted coin and gilded idol, neither overly warm nor cold
\item \textbf{Tyrian Purple} --- The historical dye drawn from murex shells, priced like gold, reserved for imperial use---deep, and arcane
\end{palettebullets}
\par\bigskip
% ---------------------------------------------------------------------------
\bookpalettecardhead{11\enspace Greece}
\palettesubsec{Colour Philosophy}
\noindent The marble is Parian, warmer and more ivory than Carrara, the actual material of the Parthenon; the democratic blue is taken from the mid-tone azure of the deep Aegean, the very sky and sea Athenians watched from the Agora. With marble as ground and democratic blue as the visual centre of gravity, gold accents, and terracotta-olive-grape as the three creative tones of narrative, wisdom and mystery.

\begin{fullwidth}\small\noindent
\begin{tabularx}{\linewidth}{@{}clllX@{}}
\toprule\rowcolor{booktableheadcolor}
\booktableheadcell{} & \booktableheadcell{Name} & \booktableheadcell{Original} & \booktableheadcell{RGB} & \booktableheadcell{Gloss} \\
\midrule
\bookcardswatch{245,240,228} & Parian Marble & \foreignname{Παριανός Μάρμαρος} \textperiodcentered\ Parianos Marmaros & 245, 240, 228 & The actual stone of the Parthenon, warmer and more ivory than Carrara \\
\bookcardswatch{212,175,55} & Gloria Aurum & \foreignname{Χρυσός} \textperiodcentered\ Chrysos & 212, 175, 55 & The same neutral gold as Rome, the highlight of gods and rite \\
\bookcardswatch{48,105,175} & Agora Kyanos & \foreignname{Κυανός Αγοράς} \textperiodcentered\ Kyanos Agoras & 48, 105, 175 & The mid-tone azure of the deep Aegean, the sky and sea Athenians watched from the Agora \\
\bookcardswatch{188,82,38} & Attic Terracotta & \foreignname{Αττικός Πηλός} \textperiodcentered\ Attikos Pelos & 188, 82, 38 & The ground of red-figure pottery, the most iconic medium colour for recording myth and epic \\
\bookcardswatch{98,128,48} & Athena's Olive & \foreignname{Ελιά Αθηνάς} \textperiodcentered\ Elia Athenas & 98, 128, 48 & The olive tree gifted by Athena to the polis, the colour of peace and wisdom \\
\bookcardswatch{90,42,92} & Dionysian Grape & \foreignname{Διονυσιακή Σταφυλή} \textperiodcentered\ Dionysiake Staphyle & 90, 42, 92 & The deep grape of wine-making, the soul colour of Dionysus, theatre, and revelry \\
\bottomrule
\end{tabularx}\end{fullwidth}

\palettesubsec{Colour Logic}
\begin{palettebullets}
\item \textbf{Parian Marble} --- Parian Marble as the ground is the material substrate of Greek architecture and sculpture---warmer and more ivory than Roman Carrara, carrying the body temperature of filtered Mediterranean sun
\item \textbf{Agora Kyanos} --- Agora Kyanos commands the visual centre of gravity---both Aegean sky and sea, and the colour of "public space" Athenian citizens saw in the Agora
\item \textbf{Gloria Aurum} --- Gloria Aurum, as accent, carries the visual weight of divinity in the smallest area---shared with Rome but used with more restraint here
\item \textbf{Attic Terracotta} --- Attic Terracotta is the "narrative colour" of Greece---red-figure pottery pressed myth and epic into this one tone, the most diagnostic medium of Greek culture
\item \textbf{Athena's Olive} --- Athena's Olive and Attic Terracotta form a warm-cool mid-balance, a colour counterpoint of wisdom and narrative
\item \textbf{Dionysian Grape} --- Dionysian Grape is the deepest anchor, the colour of Dionysus and theatre---the other pole of Greek culture: the revelry beneath the sacred
\end{palettebullets}
\par\bigskip
% ---------------------------------------------------------------------------
\bookpalettecardhead{12\enspace Kanagawa}
\palettesubsec{Colour Philosophy}
\noindent The Great Wave's revolution lay in Hokusai's introduction of newly-imported Prussian blue (Bero-ai) into ukiyo-e: a single colour family in light-and-dark variation built the drama of the entire image, while ink outline and the distant blank of Mount Fuji produced a uniquely compressed visual aesthetic. The palette centres on the Prussian blue family---from sumi-ink to deep sea, to wave-blue, to Fuji-haze, to twilight sky, to foam-white---a complete luminance ladder, like a wave rising from the depths and bursting into snow at its crest.

\begin{fullwidth}\small\noindent
\begin{tabularx}{\linewidth}{@{}clllX@{}}
\toprule\rowcolor{booktableheadcolor}
\booktableheadcell{} & \booktableheadcell{Name} & \booktableheadcell{Original} & \booktableheadcell{RGB} & \booktableheadcell{Gloss} \\
\midrule
\bookcardswatch{237,233,222} & Wave Foam & \foreignname{波白} \textperiodcentered\ Nami-shiro & 237, 233, 222 & The foam at the cresting wave, not pure white but warmed by washi paper \\
\bookcardswatch{208,224,238} & Twilight Sky & \foreignname{暮天} \textperiodcentered\ Boten & 208, 224, 238 & The horizon sky\'s gradient, Prussian blue diluted to vastness \\
\bookcardswatch{150,186,210} & Fuji Haze & \foreignname{富士霞} \textperiodcentered\ Fuji-gasumi & 150, 186, 210 & The pale-blue silhouette of distant Mount Fuji, blurred and still \\
\bookcardswatch{26,78,132} & Prussian Indigo & \foreignname{ベロ藍} \textperiodcentered\ Bero-ai & 26, 78, 132 & The main body of the image---the Prussian blue wave, the most iconic imported pigment of Edo \\
\bookcardswatch{13,38,76} & Deep Sea & \foreignname{深海} \textperiodcentered\ Shinkai & 13, 38, 76 & The ink-blue at the deepest trough of the wave, almost black yet keeping the chill of blue \\
\bookcardswatch{29,25,35} & Sumi Ink & \foreignname{墨} \textperiodcentered\ Sumi & 29, 25, 35 & The ink line of woodblock print, the skeleton of the whole image, not pure black but with a deep violet base \\
\bottomrule
\end{tabularx}\end{fullwidth}

\palettesubsec{Colour Logic}
\begin{palettebullets}
\item \textbf{Nami-shiro} --- Nami-shiro is the instant of breaking at the crest, warmed by washi paper---the only highlight in the palette
\item \textbf{Boten} --- Boten is the transitional sky at the horizon, Prussian blue diluted almost to vanishing---the colour the image breathes through
\item \textbf{Fuji-gasumi} --- Fuji-gasumi is the pale-blue distant Fuji, the single "retreat" in the painting, pushing focus back to the wave itself
\item \textbf{Bero-ai} --- Bero-ai is the wave's body, the painting's revolution---a single imported pigment carrying the drama of the entire work
\item \textbf{Shinkai} --- Shinkai is the ink-blue at the depth, the carrier of "pressure" and "abyss", almost black yet still blue
\item \textbf{Sumi} --- Sumi is the skeleton of the woodblock, governing the whole image---the "form" of the palette, without which the six blues would have nothing to attach to
\end{palettebullets}
\par\bigskip
% ---------------------------------------------------------------------------
\bookpalettecardhead{13\enspace Starry Night}
\palettesubsec{Colour Philosophy}
\noindent Van Gogh's Starry Night is not Impressionism but the summit of Post-Impressionism---he no longer catches the instant of light but burns the emotion inside directly into the brush stroke. Cobalt and ultramarine form the bone, the swirling stroke gives the night sky life, and the chrome yellow of moon and stars erupts under the weight of blue, while the cypress climbs like a black flame against the firmament; the village glow is the only warmth in the painting that belongs to the human world. Three layers of blue form an emotional axis from deepest to lightest, mimicking the layered rhythm of Van Gogh's spiral stroke.

\begin{fullwidth}\small\noindent
\begin{tabularx}{\linewidth}{@{}clllX@{}}
\toprule\rowcolor{booktableheadcolor}
\booktableheadcell{} & \booktableheadcell{Name} & \booktableheadcell{Original} & \booktableheadcell{RGB} & \booktableheadcell{Gloss} \\
\midrule
\bookcardswatch{240,208,68} & Lunar Chrome & Lumi\`ere Lunaire & 240, 208, 68 & The exploding halo of moon and starlight, Van Gogh\'s chrome yellow as the utmost tension of life \\
\bookcardswatch{198,140,52} & Village Amber & Lueurs du Village & 198, 140, 52 & The warm yellow from village windows, the only warm-cool stand-off against the cold blue sky \\
\bookcardswatch{105,155,200} & Glacial Dawn & Aube Glac\'ee & 105, 155, 200 & The pale-blue transition at the swirl\'s edge, the most delicate breath between deep night and moonlight \\
\bookcardswatch{48,96,165} & Swirling Ultramarine & Tourbillon Outremer & 48, 96, 165 & The ultramarine of the swirl\'s body, the signature Post-Impressionist colour---not natural but emotional \\
\bookcardswatch{22,50,30} & Cypress Nocturne & Cypr\`es Nocturne & 22, 50, 30 & The deep dark green of the foreground cypress, nearly black, the only thing in the painting that "refuses the light" \\
\bookcardswatch{20,36,88} & Midnight Cobalt & Minuit Cobalt & 20, 36, 88 & The deepest base of the night sky, dense with oppressive emotional weight---the abyss of Van Gogh\'s inner world \\
\bottomrule
\end{tabularx}\end{fullwidth}

\palettesubsec{Colour Logic}
\begin{palettebullets}
\item \textbf{Lumi\`ere Lunaire} --- Lumière Lunaire is the painting's "burst point", moon and stars exploding under the weight of blue---the very limit of life's tension
\item \textbf{Lueurs du Village} --- Lueurs du Village is the only earthly warmth in the painting, standing against the cold blue sky as the sole warm-cool counterpoint
\item \textbf{Aube Glac\'ee} --- Aube Glacée is the transition, the most delicate breath between deep night and moonlight, between oppression and release
\item \textbf{Tourbillon Outremer} --- Tourbillon Outremer is the body of the swirl, the signature Post-Impressionist hue---not natural but emotional, the "melody line" of the palette
\item \textbf{Cypr\`es Nocturne} --- Cyprès Nocturne is the only element that "refuses the light", almost black---the visual anchor, holding up the contrast of the whole work with the darkest tone
\item \textbf{Minuit Cobalt} --- Minuit Cobalt is the deepest emotional base of the night sky, Van Gogh's inner abyss, from which every other blue grows
\end{palettebullets}
\par\bigskip
% ---------------------------------------------------------------------------
\bookpalettecardhead{14\enspace A Thousand Li}
\palettesubsec{Colour Philosophy}
\noindent A Thousand Li of Rivers and Mountains was painted by the Northern-Song prodigy Wang Ximeng with mineral pigments layered repeatedly---azurite and malachite ground fine and pressed coat upon coat, the "thin yet thick" technique sinking colour into the silk to give it a mineral translucence that glows from within outward. The six colours do not stand alone; together they form a colour ecology in which mountain and water breathe as one.

\begin{fullwidth}\small\noindent
\begin{tabularx}{\linewidth}{@{}clllX@{}}
\toprule\rowcolor{booktableheadcolor}
\booktableheadcell{} & \booktableheadcell{Name} & \booktableheadcell{Original} & \booktableheadcell{RGB} & \booktableheadcell{Gloss} \\
\midrule
\bookcardswatch{218,203,170} & Song Silk & \foreignname{宋绢} \textperiodcentered\ S\`ong-ju\=an & 218, 203, 170 & The warm rice ground of millennia-aged silk, not pigment but the colour of time itself \\
\bookcardswatch{110,165,195} & Sky Azurite & \foreignname{烟渚} \textperiodcentered\ Y\=an-zh\v{u} & 110, 165, 195 & The diluted layer of azurite and lead white, used for distant mountain silhouette and river highlights \\
\bookcardswatch{183,118,45} & Ochre-Gold & \foreignname{砂赭} \textperiodcentered\ Sh\=a-zh\v{e} & 183, 118, 45 & Ochre limonite mixed with mud-gold, used for reef texture, dwelling, gallery bridge, and gold lines on water \\
\bookcardswatch{96,148,110} & Pale Malachite & \foreignname{岚霭绿} \textperiodcentered\ L\'an-\v{a}i-l\"u & 96, 148, 110 & The transitional tone of progressively diluted malachite, the mid-ground foliage and the cloud-mountain interface \\
\bookcardswatch{48,105,76} & Malachite True & \foreignname{苍岭绿} \textperiodcentered\ C\=ang-l\v{i}ng-l\"u & 48, 105, 76 & The true colour of natural malachite, layered on sunlit mountain faces, the most life-bearing tone of the painting \\
\bookcardswatch{35,78,112} & Azurite Deep & \foreignname{霁山青} \textperiodcentered\ J\`i-sh\=an-q\=ing & 35, 78, 112 & The deep mineral blue of finely-ground azurite, the dominant tone of the mountain\'s shaded face and deep water \\
\bottomrule
\end{tabularx}\end{fullwidth}

\palettesubsec{Colour Mapping of the Five Characteristics}
\begin{palettebullets}
\item \textbf{Azure-and-Green Lead} --- Azurite Deep and Malachite True jointly anchor the main key---one cool one warm, one yin one yang, each balancing rather than competing with the other
\item \textbf{Gradient Layering} --- from Azurite Deep to Sky Azurite, from Malachite True to Pale Malachite---each line is a luminance ladder, mimicking the layered technique of "thin yet thick"
\item \textbf{Gold-Ochre Accent} --- Ochre-Gold takes on the largest duty of temperature adjustment in the smallest area, keeping the azure-green lead from slipping into chill
\item \textbf{Ethereal Refinement} --- Sky Azurite and Pale Malachite open breath between the dense mineral tones, so the whole palette neither clogs nor stifles
\item \textbf{Antique Silk Quality} --- Song Silk holds the five tones, the way millennia-aged silk settles every intensity into a far, restrained tone
\end{palettebullets}
\par\bigskip
% ---------------------------------------------------------------------------
\bookpalettecardhead{15\enspace And Quiet Flows the Don}
\palettesubsec{Colour Philosophy}
\noindent Sholokhov spent twenty years writing this epic. The fate of the Don Cossacks never lives in the bright places---it is buried in frozen earth, rusted onto sabres, congealed in the black blood at the trench, scattered through reed marshes the snow has crossed. This palette refuses every "colour" and keeps only what time leaves after wearing them down: not blue but turbid blue, not red but dark blood, not yellow but withered grass, not grey but lead grey.

\begin{fullwidth}\small\noindent
\begin{tabularx}{\linewidth}{@{}clllX@{}}
\toprule\rowcolor{booktableheadcolor}
\booktableheadcell{} & \booktableheadcell{Name} & \booktableheadcell{Original} & \booktableheadcell{RGB} & \booktableheadcell{Gloss} \\
\midrule
\bookcardswatch{142,120,70} & Wormwood Yellow & \foreignname{Полынь} \textperiodcentered\ Polyn & 142, 120, 70 & Reeds and wormwood by the Don after the first frost, a withered yellow-green \\
\bookcardswatch{110,108,104} & Steppe Lead Grey & \foreignname{Степной Пепел} \textperiodcentered\ Stepnoy Pepel & 110, 108, 104 & The winter steppe where sky and ground merge into one tone, the very grey of a Cossack man\'s gaze \\
\bookcardswatch{118,86,52} & Don Earth & \foreignname{Донская Земля} \textperiodcentered\ Donskaya Zemlya & 118, 86, 52 & The chestnut soil unique to the Don basin, brown with the smell of humus \\
\bookcardswatch{122,66,36} & Rusted Iron & \foreignname{Ржавое Железо} \textperiodcentered\ Rzhavoye Zhelezo & 122, 66, 36 & The rust on sabre, stirrup, and rifle stock---the colour of war corroding metal \\
\bookcardswatch{52,70,90} & Don Muddy Blue & \foreignname{Мутный Дон} \textperiodcentered\ Mutnyy Don & 52, 70, 90 & The Don is never clear---it carries silt and the blood of the slain, the time-base of the whole novel \\
\bookcardswatch{88,26,24} & Caked Blood & \foreignname{Запёкшаяся Кровь} \textperiodcentered\ Zapyokshayasya Krov & 88, 26, 24 & The deep brown-red that seeps into uniform three days after a wound---blood as fate, silenced \\
\bottomrule
\end{tabularx}\end{fullwidth}

\palettesubsec{Tragic Logic of the Colour Structure}
\begin{palettebullets}
\item \textbf{Cold Grey Holds the Lid} --- Steppe Lead Grey and Don Muddy Blue press every warm tone beneath them, as history's grinding of individual fate has never let go
\item \textbf{Earth Tones Bear the Weight} --- Don Earth and Wormwood Yellow are the only "living" still left to this land, but desaturated and dulled---even beauty is tired beauty
\item \textbf{Dark Blood Anchors the Floor} --- Caked Blood is the deepest extremity, never the protagonist, appearing only in the heaviest shadow---like death in this novel, never dramatised
\item \textbf{Rust Threads Through} --- Rusted Iron crosses between earth-brown and dark blood, the carrier of "time-sense"---telling the viewer that all of this has already been long past
\end{palettebullets}
\par\bigskip
% ---------------------------------------------------------------------------
\bookpalettecardhead{16\enspace Cyberpunk Edgerunners}
\palettesubsec{Colour Philosophy}
\noindent The visual language of Edgerunners is overload aesthetics in essence---colour here is not ornament but symptom. David's breakdown comes with the runaway erosion of fluorescent tones; the neon of Night City is not the sign of prosperity but an anaesthetic; Lucy's blue is the only real exit this world owns. The central contradiction of the palette is this: the darkest possible base set against the most fluorescent possible highlight, producing a vibration close to bodily discomfort.

\begin{fullwidth}\small\noindent
\begin{tabularx}{\linewidth}{@{}clllX@{}}
\toprule\rowcolor{booktableheadcolor}
\booktableheadcell{} & \booktableheadcell{Name} & \booktableheadcell{Original} & \booktableheadcell{RGB} & \booktableheadcell{Gloss} \\
\midrule
\bookcardswatch{225,255,8} & Psycho Yellow & \foreignname{サイコ黄} \textperiodcentered\ Saiko-ki & 225, 255, 8 & The soul colour of David\'s jacket, the visual overload at the onset of cyberpsychosis \\
\bookcardswatch{10,238,100} & Flatline Green & \foreignname{電脳緑} \textperiodcentered\ Denn\=o{}-ryoku & 10, 238, 100 & The code stream of network intrusion and the subdermal light tracery of chromed bodies \\
\bookcardswatch{238,18,120} & Neon Magenta & \foreignname{ネオン桃} \textperiodcentered\ Neon-momo & 238, 18, 120 & The most overflowing neon of Night City, the colour of prosperity\'s lie \\
\bookcardswatch{205,20,35} & Edgerunner Red & \foreignname{エッジ赤} \textperiodcentered\ Edge-aka & 205, 20, 35 & The price an edgerunner pays, the colour on the concrete after a street firefight \\
\bookcardswatch{15,32,98} & Luna Blue & \foreignname{月の蒼} \textperiodcentered\ Tsuki-no-ao & 15, 32, 98 & Lucy looking up at the moon---the only colour that does not belong to Night City \\
\bookcardswatch{14,8,28} & Night City Void & \foreignname{夜都虚空} \textperiodcentered\ Yato Kok\=u{} & 14, 8, 28 & The absolute darkness at the floor of Night City\'s street, carrying a deep violet undertone \\
\bottomrule
\end{tabularx}\end{fullwidth}

\palettesubsec{Narrative Logic of the Colour Structure}
\begin{palettebullets}
\item \textbf{Void Holds the Floor} --- Night City Void is the absolute ground, with saturation near zero, the precondition that lets the other five "explode". Without darkness, neon is only pigment
\item \textbf{Cold/Warm Rupture} --- Luna Blue against Neon Magenta and Psycho Yellow forms the most extreme warm-cool opposition---blue is Lucy's lunar dream; pink and yellow are the lies and madness of Night City
\item \textbf{Fluorescent Overload} --- Psycho Yellow and Flatline Green are the visual core of cyberpsychosis aesthetics; placed side by side they produce a vibration close to bodily discomfort
\item \textbf{Red Finale} --- Edgerunner Red is not the protagonist in the plate, but the most honest colour---like the ending, impossible to avoid
\end{palettebullets}
\par\bigskip
% ---------------------------------------------------------------------------
\bookpalettecardhead{17\enspace The Grand Budapest Hotel}
\palettesubsec{Colour Philosophy}
\noindent Wes Anderson's colour is never realist---it is a precisely calibrated emotional filter. The pink of The Grand Budapest Hotel is not pink; it is the colour of an entire vanishing old-European aristocratic world compressed into a single tone. The purple is not purple; it is Monsieur Gustave's final, never-impolite dignity. Every shade has been deliberately de-realised.

\begin{fullwidth}\small\noindent
\begin{tabularx}{\linewidth}{@{}clllX@{}}
\toprule\rowcolor{booktableheadcolor}
\booktableheadcell{} & \booktableheadcell{Name} & \booktableheadcell{Original} & \booktableheadcell{RGB} & \booktableheadcell{Gloss} \\
\midrule
\bookcardswatch{247,235,215} & Vanilla Cream & Cr\`eme Vanille & 247, 235, 215 & Hotel interior, Mendl\'s sponge, letter paper---the very air of a fairy-tale world \\
\bookcardswatch{168,178,196} & Alpine Mist & Brume Alpine & 168, 178, 196 & Snowy alps, the desaturated cool outside a winter coach---the emotional refrigerant \\
\bookcardswatch{237,148,158} & Mendl's Rose & Rose M\'endl & 237, 148, 158 & The signature pink of the hotel façade and Mendl\'s confectionery box, a dust-touched old rose \\
\bookcardswatch{176,136,60} & Antique Gold & Dor\'e Antique & 176, 136, 60 & Gilded mirror frame, doorknob, shoulder braid---the patina of old aristocracy \\
\bookcardswatch{128,88,148} & Concierge Violet & Violet Concierge & 128, 88, 148 & The purple of Monsieur Gustave\'s uniform, at once authority and absurdity \\
\bookcardswatch{143,52,65} & Vintage Bordeaux & Bordeaux Vintage & 143, 52, 65 & The dramatic anchor of dark scenes---library leather, claret, conspiracy and danger \\
\bottomrule
\end{tabularx}\end{fullwidth}

\palettesubsec{Colour Mapping of the Five Aesthetic Rules}
\begin{palettebullets}
\item \textbf{Macaron Softness} --- Rose Méndl and Violet Concierge carry the main visual, high saturation deliberately tamed in luminance---like the shell of a macaron, distinct yet never glaring
\item \textbf{Cream Warm Ground} --- Crème Vanille rules through the largest negative space, the "void" of Wes Anderson's symmetric composition---keeping the other five from crowding
\item \textbf{Antique Dark Gold} --- Doré Antique carries the weight of "aristocracy", taking no part in the fairy-tale narrative---glowing only in frame and trim
\item \textbf{Cool Restraint} --- Brume Alpine is the only cool, its low saturation making it the regulator, opening a window of breath inside an otherwise sweet warm field
\item \textbf{Sweet Without Cloy} --- Bordeaux Vintage is the most subtle move---a deep dark anchor, the drop of amaretto inside the dessert, giving every lightness its weight
\end{palettebullets}
\par\bigskip
% ---------------------------------------------------------------------------
\bookpalettecardhead{18\enspace Renaissance Florence}
\palettesubsec{Colour Philosophy}
\noindent The colour of Florence was not designed but grew out of substance itself---lapis lazuli ground into pigment, lead white folded with yolk into tempera ground, gold leaf laid on wooden altar, terracotta fired in the Tuscan kiln into dome tile. All six of these colours have a real material provenance: they were once pigment, ore, building stone, and timber---the spectrum Botticelli, Giotto, and Donatello breathed together.

\begin{fullwidth}\small\noindent
\begin{tabularx}{\linewidth}{@{}clllX@{}}
\toprule\rowcolor{booktableheadcolor}
\booktableheadcell{} & \booktableheadcell{Name} & \booktableheadcell{Original} & \booktableheadcell{RGB} & \booktableheadcell{Gloss} \\
\midrule
\bookcardswatch{238,228,208} & Florentine Ivory & Avorio Fiorentino & 238, 228, 208 & The lead-white ground of cut Carrara marble and tempera panel, warmed by yolk and linseed \\
\bookcardswatch{190,148,52} & Altar Gold & Oro dell'Altare & 190, 148, 52 & The gold leaf of the Madonna\'s nimbus and Medici ornament, an aged gold still glowing after five centuries of oxidation \\
\bookcardswatch{172,84,50} & Brunelleschi Cotto & Cotto Brunellesco & 172, 84, 50 & The signature terracotta of Brunelleschi\'s dome and Florentine eaves \\
\bookcardswatch{54,80,60} & Cypress Verde & Verde Cipresso & 54, 80, 60 & The cypress and verdigris mix of the Tuscan hill \\
\bookcardswatch{46,70,118} & Lapis Ultramarine & Oltremare di Lapislazzuli & 46, 70, 118 & Ground from Afghan lapis lazuli, priced like gold, reserved for the Madonna\'s robe \\
\bookcardswatch{85,54,34} & Tuscan Walnut & Noce Toscano & 85, 54, 34 & The walnut panel and dark underglaze of the Florentine workshop \\
\bottomrule
\end{tabularx}\end{fullwidth}

\palettesubsec{Colour Mapping of the Four Temperaments}
\begin{palettebullets}
\item \textbf{Sacred Solemnity} --- Oltremare carries the spiritual weight of the whole set. In the Renaissance, lapis was priced like gold and not available to common commission---its very presence is a sacred declaration; Oro dell'Altare lets out its glow in the most restrained way, and the two together build the rite of religious art
\item \textbf{Warm Nobility} --- Avorio Fiorentino and Oro dell'Altare are the two faces of the aristocratic ground---one the warm touch of marble and silk, the other the Medici memory of wealth feeding art; Cotto Brunellesco then tempers sacred austerity with human warmth
\item \textbf{Classical Depth} --- All six tones are deliberately de-modernised, no pure colour anywhere, each carrying the "impurity" of mineral debris---this is the most fundamental difference between natural pigment and synthetic dye. Heft comes from restrained saturation, depth from each colour's inner complexity
\item \textbf{Humanist Refinement} --- Verde Cipresso and Noce Toscano are the humanist anchors of the palette---the green comes from the natural Tuscan hill, the brown from the walnut panel of the workshop; together they pull the sacred from sky back to earth
\end{palettebullets}
\par\bigskip
% ---------------------------------------------------------------------------
\bookpalettecardhead{19\enspace Soviet Avant-Garde}
\palettesubsec{Colour Philosophy}
\noindent The Soviet avant-garde never "designed" colour---they declared it. Rodchenko said colour must have a task; Lissitzky said form is politics; Malevich said the Suprematist black square is the absolute of zero. This palette refuses harmony, refuses gradient, refuses the decorative transition: every colour is a hard-edged manifesto. Red does not explain itself, black does not apologise, grey does not curry favour with anyone.

\begin{fullwidth}\small\noindent
\begin{tabularx}{\linewidth}{@{}clllX@{}}
\toprule\rowcolor{booktableheadcolor}
\booktableheadcell{} & \booktableheadcell{Name} & \booktableheadcell{Original} & \booktableheadcell{RGB} & \booktableheadcell{Gloss} \\
\midrule
\bookcardswatch{225,218,205} & Newsprint White & \foreignname{ГАЗЕТА} \textperiodcentered\ Gazeta & 225, 218, 205 & The newspaper white of Pravda and constructivist brochures, an ink-soaked warm grey \\
\bookcardswatch{200,150,32} & Propaganda Ochre & \foreignname{ПЛАКАТ} \textperiodcentered\ Plakat & 200, 150, 32 & The ochre of Klutsis photomontage poster and collectivisation propaganda print \\
\bookcardswatch{118,114,108} & Concrete Grey & \foreignname{БЕТОН} \textperiodcentered\ Beton & 118, 114, 108 & The poured concrete of Tatlin\'s Tower, Moscow factory, and Five-Year-Plan worksite \\
\bookcardswatch{50,80,138} & Blueprint Blue & \foreignname{ЧЕРТЁЖ} \textperiodcentered\ Chertyozh & 50, 80, 138 & The blueprint blue of the working-class uniform, deliberately desaturated \\
\bookcardswatch{196,28,28} & Lenin Red & \foreignname{КРАСНЫЙ} \textperiodcentered\ Krasnyy & 196, 28, 28 & The revolutionary red of the Bolshevik flag and Rodchenko poster, pressed down in saturation to weigh like iron \\
\bookcardswatch{24,20,18} & Cast-Iron Black & \foreignname{ЧЁРНЫЙ} \textperiodcentered\ Chyornyy & 24, 20, 18 & The absolute black of Malevich\'s Black Square, where nothingness and certainty coexist \\
\bottomrule
\end{tabularx}\end{fullwidth}

\palettesubsec{Colour Mapping of the Six Aesthetic Rules}
\begin{palettebullets}
\item \textbf{Red-Black Hard Strike} --- \foreignname{КРАСНЫЙ} and \foreignname{ЧЁРНЫЙ} are the structural core, with no transition and no buffer between them---contrast itself is the argument
\item \textbf{Industrial Grey Axis} --- \foreignname{БЕТОН} takes part in no drama, it merely stands and bears the load---like the nameless workers of the Five-Year-Plan, the spine, not the focus
\item \textbf{Mechanical Blue Precision} --- \foreignname{ЧЕРТЁЖ} is pressed nearly to colourlessness, forming an ideological warm-cool opposition with \foreignname{КРАСНЫЙ}---red the passion of revolution, blue the discipline of construction
\item \textbf{Democratic Coarse White} --- \foreignname{ГАЗЕТА} is the only "anti-elite" tone, refusing the aristocratic warmth of ivory and choosing the coarseness of cheap print
\item \textbf{Unkept Ochre} --- \foreignname{ПЛАКАТ} is the most tragic tone, the colour of poster wheat and forever-sunlit utopian tomorrows---a colour of promise, not of reality
\end{palettebullets}
\par\bigskip
% ---------------------------------------------------------------------------
\bookpalettecardhead{20\enspace Constantinople}
\palettesubsec{Colour Philosophy}
\noindent Constantinople is the only city in human history to belong, all at once, to two continents, three civilisations, and four religions---its colours are never simple. The gold of Hagia Sophia's dome holds Greek craft, Roman engineering, and Eastern theology in superposition. The deep blue of the Bosporus mixes the Black Sea's chill with the Aegean's warmth. The imperial purple is luxury distilled from Mediterranean murex.

\begin{fullwidth}\small\noindent
\begin{tabularx}{\linewidth}{@{}clllX@{}}
\toprule\rowcolor{booktableheadcolor}
\booktableheadcell{} & \booktableheadcell{Name} & \booktableheadcell{Original} & \booktableheadcell{RGB} & \booktableheadcell{Gloss} \\
\midrule
\bookcardswatch{236,225,207} & Proconnesian Ivory & \foreignname{Λευκός} \textperiodcentered\ Fildi\c{s}i & 236, 225, 207 & Marble of the Sea of Marmara, the floor of Hagia Sophia \\
\bookcardswatch{195,150,42} & Hagia Sophia Gold & \foreignname{Χρυσός} \textperiodcentered\ Alt\i n & 195, 150, 42 & The aged gold smoked by fifteen centuries of candles, sacred and unshowy \\
\bookcardswatch{170,112,50} & Theodosian Ochre & \foreignname{\'{Ω}χρα} \textperiodcentered\ Toprak & 170, 112, 50 & The brick of the Theodosian Walls, the shared ground of Greek, Roman, Anatolian \\
\bookcardswatch{50,76,58} & Cypress Dark Green & \foreignname{Κυπαρίσσι} \textperiodcentered\ Selvi & 50, 76, 58 & Byzantine cemetery and Islamic garden---two faiths watching the same eternity \\
\bookcardswatch{102,32,65} & Imperial Murex & \foreignname{Πορφύρα} \textperiodcentered\ Mor & 102, 32, 65 & Restricted to the Byzantine emperor: a single gram demands ten thousand murex shells \\
\bookcardswatch{30,52,88} & Bosporus Blue & \foreignname{Βόσπορος} \textperiodcentered\ Bo\u{g}az & 30, 52, 88 & The deep blue of the strait at dusk, where the Black Sea meets the Aegean \\
\bottomrule
\end{tabularx}\end{fullwidth}

\palettesubsec{The Colour Structure of Civilisations in Confluence}
\begin{palettebullets}
\item \textbf{East-West Warm/Cool Axis} --- Bosporus blue (Western reason, maritime order) and Theodosian ochre (Eastern earth, Anatolian soil) form the East-West civilisational axis---opposite in temperature, they coexisted for a millennium without dissolving in this city
\item \textbf{Sacred Triangle} --- Hagia Sophia gold, imperial purple, Proconnesian ivory together build the colour space of the sacred---gold is the light of God, purple is God's human agent, ivory is the architectural skin in which God resides
\item \textbf{Purple as Civilisational Bridge} --- Imperial Murex bears the most complex role in the whole set, belonging at once to Roman imperial tradition, Byzantine Christian theology, and Persian-Eastern court luxury
\item \textbf{Dark Green Holds Deep Time} --- Cypress dark green carries the longest time span in the most low-key form---present at once in Byzantine cemetery cypress, Ottoman garden, and the verdigris dome of the mosque
\item \textbf{Restrained Gold} --- The palette deliberately compresses Hagia Sophia gold into mineral-tone saturation, glowing without glare---this is the central logic of Byzantine mosaic: gold is not ornament but the materialisation of light
\end{palettebullets}
\par\bigskip
% ---------------------------------------------------------------------------
\bookpalettecardhead{21\enspace France}
\palettesubsec{Colour Philosophy}
\noindent French colour is never simply colour---it is the material appearance of a worldview. The Republican red holds the guillotine and the rose; the deep blue houses both Descartes' reason and Monet's morning mist; champagne ivory is the same ground tone Versailles marble shares with Left-Bank café paper. The palette tries to catch the central inner contradiction of the French spirit: extreme reason and extreme romance, coexisting on the same patch of land.

\begin{fullwidth}\small\noindent
\begin{tabularx}{\linewidth}{@{}clllX@{}}
\toprule\rowcolor{booktableheadcolor}
\booktableheadcell{} & \booktableheadcell{Name} & \booktableheadcell{Original} & \booktableheadcell{RGB} & \booktableheadcell{Gloss} \\
\midrule
\bookcardswatch{180,32,40} & Marianne Red & Rouge Marianne & 180, 32, 40 & The Tricolour red, but no flat red---a mid-value between Bordeaux\'s depth and the passion of the Phrygian cap \\
\bookcardswatch{232,222,205} & Champagne Ivory & Ivoire Champagne & 232, 222, 205 & Not pure white---champagne effervescence, Louvre limestone, and Left-Bank café stationery \\
\bookcardswatch{44,74,138} & Republic Blue & Bleu R\'epublique & 44, 74, 138 & The Tricolour blue, the royal blue of Sèvres porcelain---the blue of Enlightenment reason \\
\bookcardswatch{143,108,148} & Provence Lavender & Lavande Proven\c{c}ale & 143, 108, 148 & The true colour of a Provençal lavender field between afternoon sun and chalk-grey soil \\
\bookcardswatch{142,133,58} & Versailles Olive & Olive Dor\'ee & 142, 133, 58 & The gold-green of clipped box hedge in the formal garden of Versailles---the colour of classical French parterre aesthetics \\
\bookcardswatch{122,126,132} & Parisian Zinc & Zinc Parisien & 122, 126, 132 & The reflection of Haussmannian Parisian zinc roof under overcast sky \\
\bottomrule
\end{tabularx}\end{fullwidth}

\palettesubsec{Colour Structure of the French Spirit}
\begin{palettebullets}
\item \textbf{Tricolour Reread} --- Rouge Marianne, Ivoire Champagne, and Bleu République are not the flag reproduced but its spirit de-symbolised and re-distilled---red is given the weight of Bordeaux, white the temperature of time, blue the depth of the Enlightenment. The flag is political; this palette is cultural
\item \textbf{North-South Tension} --- Lavande Provençale and Zinc Parisien form France's deepest spiritual opposition---the south is the sense, soil, sun, and lazy romance; Paris is reason, geometry, grey sky, and restrained elegance
\item \textbf{Avant-Garde Strike Latent} --- Olive Dorée placed against Bleu République is the most avant-garde pairing in the set---one belongs to the Sun King's classical order, the other to the rational edge of republican revolution---hues close, semantics opposed
\item \textbf{Restrained Elegance} --- Every tone is deliberately tamed in saturation---red does not burn, blue does not chill, violet does not flirt, gold does not boast, grey does not deaden, white does not glare. The French sense of "elegance" has never been the intensity of colour, but its trained, exact self-restraint
\item \textbf{Earthbound Warmth} --- Ivoire Champagne and Lavande Provençale weave the palette's "earth-sense" baseline---one the warmth of limestone and wheat, the other the warmth of dried lavender and chalk soil. French romance has never hung in midair
\end{palettebullets}
\par\bigskip
% ---------------------------------------------------------------------------
\bookpalettecardhead{22\enspace Kyoto}
\palettesubsec{Colour Philosophy}
\noindent The colours of Kyoto never speak first---they wait to be seen. The moss green of the Zen courtyard does not compete, the white gravel of the karesansui says nothing, the tea-brown of machiya cedar is silent in the rain, the deep-indigo night sky presses the thousand-year capital without a sound. Only the vermilion torii, amid all this restraint, burns the highest temperature in the smallest area---like a single stem in the tea room, the one point of heat inside the cold and still.

\begin{fullwidth}\small\noindent
\begin{tabularx}{\linewidth}{@{}clllX@{}}
\toprule\rowcolor{booktableheadcolor}
\booktableheadcell{} & \booktableheadcell{Name} & \booktableheadcell{Original} & \booktableheadcell{RGB} & \booktableheadcell{Gloss} \\
\midrule
\bookcardswatch{234,226,212} & Silk White & \foreignname{絹白} \textperiodcentered\ Kinushiro & 234, 226, 212 & The ground of Nishijin silk and washi paper, the natural warm tone of silkworm and paper-mulberry \\
\bookcardswatch{110,84,64} & Withered Tea & \foreignname{枯茶} \textperiodcentered\ Karacha & 110, 84, 64 & The brown of a century-old machiya cedar pillar shaped by time and hearth smoke---the colour of Riky\=u{}\'s wabi \\
\bookcardswatch{204,184,178} & Cherry Grey & \foreignname{桜鼠} \textperiodcentered\ Sakura-nezumi & 204, 184, 178 & Fallen cherry on a stone lantern, blurred by rain---the precise definition of mono no aware \\
\bookcardswatch{116,126,120} & Bamboo Grey & \foreignname{青竹鼠} \textperiodcentered\ Aotake-nezumi & 116, 126, 120 & The bamboo node in the dawn mist of the Arashiyama grove, a Zen statement of warm-cool balance \\
\bookcardswatch{80,94,68} & Moss Colour & \foreignname{苔色} \textperiodcentered\ Koke-iro & 80, 94, 68 & The deep moss green at the foot of the Ry\=o{}an-ji stones, the dull print time has left \\
\bookcardswatch{34,40,66} & Deep Indigo & \foreignname{紺} \textperiodcentered\ Kon & 34, 40, 66 & The deep indigo of Noh costume and Kyoto-Yuzen night patterns---the colour of meditation itself \\
\bottomrule
\end{tabularx}\end{fullwidth}

\palettesubsec{Colour Structure of the Kyoto Spirit}
\begin{palettebullets}
\item \textbf{Negative Space as Bone} --- Kinushiro and Sakura-nezumi build the palette's breathing room---the former is void, negative form, the tea-room wall; the latter is the perceptual edge of fading-not-yet-gone. Kyoto aesthetics is not about filling in, but about knowing where there should be nothing
\item \textbf{Warm/Cool Balance} --- Among the six grounds, Karacha and Koke-iro are warm, Kon and Aotake-nezumi are cool, Sakura-nezumi and Kinushiro drift in between---four-season balance, the cold stillness of the Zen courtyard exactly counterpointed with the warmth of machiya
\item \textbf{Y\=u{}gen via Low Saturation} --- All colours are deliberately stripped of brightness, mimicking Kyoto's basin morning mist---the mist does not erase colour but draws it inward, which is the very physics of y\=u{}gen
\item \textbf{Ink-Wash Eastern Structure} --- The whole palette can be read as one ink-wash---Kon is the deep ink, Koke-iro the middle ink, Karacha the light ink, Aotake-nezumi the spent ink, Sakura-nezumi the ink's diffusion, Kinushiro the paper---and an unseen vermilion is the painter's final cinnabar seal
\end{palettebullets}
\par\bigskip
% ---------------------------------------------------------------------------
\bookpalettecardhead{23\enspace Siamese Dream}
\palettesubsec{Colour Philosophy}
\noindent The cover of \emph{Siamese Dream} is itself a palette manifesto---two girls leaning into each other inside a softly out-of-focus warm halo, a faint deep green pressing in from behind, an orange light spilling through like some afternoon out of memory, beautiful past the point of belief. This is not the high-contrast assault that rock-album sleeves usually reach for; it is a deliberately constructed dream texture---overexposed, soft, tender film-grain distortion---just as Corgan poured the rawest emotion into the most distorted wall of guitars.

\begin{fullwidth}\small\noindent
\begin{tabularx}{\linewidth}{@{}clllX@{}}
\toprule\rowcolor{booktableheadcolor}
\booktableheadcell{} & \booktableheadcell{Name} & \booktableheadcell{Original} & \booktableheadcell{RGB} & \booktableheadcell{Gloss} \\
\midrule
\bookcardswatch{238,226,208} & Luna & Luna & 238, 226, 208 & The warm white halo of the cover, the lingering breath of the album's closing track---not pure white, but the trace of human warmth that survives film overexposure, the last soft light before waking \\
\bookcardswatch{182,180,193} & Disarm & Disarm & 182, 180, 193 & A washed-out grey-violet. \emph{Disarm} is the album's least-defended moment, strings and bells stripping every layer of distortion---the colour just before glass shatters: transparent, brittle, carrying the chill of childhood wounds \\
\bookcardswatch{198,158,65} & Today & Today & 198, 158, 65 & A dusted, antiqued gold. \emph{Today} sounds like the brightest melody on the record yet was born of Corgan's darkest days---so this gold must be faded, slightly bitter beneath its sweetness, beautiful and not to be trusted \\
\bookcardswatch{150,145,135} & Hummer & Hummer & 150, 145, 135 & Smoke-ash brown. \emph{Hummer} stacks wall after wall of guitar to manufacture a weightless suspension---this is the colour seen by a sleepwalker, neither awake nor asleep, the most honest material form of shoegaze aesthetics \\
\bookcardswatch{188,105,40} & Mayonaise & Mayonaise & 188, 105, 40 & A deep, settled burnt orange, the colour of the cover's light source. \emph{Mayonaise} is the slowest, most painful, most private cut on the record---this orange is not warmth but the slow, low-temperature caramelisation of an emotion held in too long \\
\bookcardswatch{42,68,48} & Soma & Soma & 42, 68, 48 & The deep green half-hidden behind the cover, the nine-minute epic that climbs from tender hush to collapse---the colour of organic matter growing in darkness, moss and rotting leaves, where beauty and suffocation share the same body \\
\bottomrule
\end{tabularx}\end{fullwidth}

\palettesubsec{The Paradox Structure of Adolescent Feeling}
\begin{palettebullets}
\item \textbf{Dream-Wrap and Distortion} --- Luna's warm white and Hummer's smoke-ash together form the record's \emph{out-of-focus layer}---every pain is first overexposed, softened, and wrapped in a wall of sound before it is allowed to be heard. Not avoidance, but releasing pain at a survivable density
\item \textbf{Beautiful Vessels for the Heaviest Contents} --- Today's tarnished gold and Disarm's grey-violet sit at opposite poles of the paradox---the former dresses self-destructive impulse in a bright melody, the latter tells a childhood wound through music-box bells. Corgan always pours the heaviest contents into the most beautiful vessel; this is the album's core emotional strategy
\item \textbf{Low-Temperature Caramelisation at the Emotional Core} --- Mayonaise's burnt orange is the palette's emotional core (and conveniently the pal5=primary slot)---not red on fire, but orange that has already burned, slow-roasted by time, weighted down by its own gravity. The heat here is not from explosion but from refusing to go out
\item \textbf{The Abyss as Deepest Anchor} --- Soma's deep green (pal6=accent) is the only \emph{black-substitute} among the six---not pure black but organic black, the black of moss, leaf-rot, and deep forest. The palette never anchors on pure black; it presses the record down with a darkness that is still alive and still growing
\item \textbf{Permeation over Saturation} --- All six colours retreat from brightness because \emph{Siamese Dream}'s feelings never explode---they seep. No single colour fights for attention; each lingers at its lowest posture, the way a shoegaze guitar sustains long after the strike
\end{palettebullets}


\end{document}
